\ioaafig{2023_TH_Q02-s1.pdf}
The ratio of the orbits is 9. According to Kepler's law, the ratio of the squares of the periods is equal
to the ratio of the cubes of the radii of the orbits,
\[
\left(\frac{r_1}{r_2}\right)^3=\left(\frac{P_1}{P_2}\right)^2.
\]
\hfill\mbox{(3 points)}

Substituting the numbers, we obtain the period of planet 2,
\[
P_2=P_1\left(\frac{r_2}{r_1}\right)^{3/2}=10\times3^3=270\ \mathrm{days}.
\]
\hfill\mbox{(2 points)}

The angular velocity is $\omega=2\pi/P$ and the orbital phase is $\varphi=\omega t$.
A transit will occur when $\varphi_1-\varphi_2=2k\pi$, with $k$ an integer. The first one will occur for $k=1$.
Therefore,
\[
2\pi t\left(\frac{1}{10}-\frac{1}{270}\right)=2\pi,
\qquad t=10.3846\ \mathrm{days}.
\]
This is 10 days, 9 hours and 14 minutes after the original transit, thus on 6 October 2023, 7:14 pm.
\hfill\mbox{(5 points)}
